\begin{pastexam}[label={px:ncku-114-cce-ds-15}]{成大電通所 114 學年度 單選題第 15 題}
（單選；每題 5 分，答錯倒扣 2 分）

Please give the maximum flow from S to T in the network flow graph in Fig 2, where the value associated to an
edge is its capacity.

\begin{center}
\begin{tikzpicture}[v/.style={circle,draw,minimum size=7.5mm},w/.style={inner sep=2pt},>={Stealth[length=2.4mm]}]
\node[v] (S) at (0,0) {};
\node[v] (a) at (2.3,1.3) {};
\node[v] (b) at (5.0,1.3) {};
\node[v] (c) at (2.3,-1.2) {};
\node[v] (d) at (5.0,-1.2) {};
\node[v] (T) at (7.3,0) {};
\draw[->] (S) -- node[w,above left] {4} (a);
\draw[->] (S) -- node[w,below left] {6} (c);
\draw[->] (a) -- node[w,above] {3} (b);
\draw[->] (b) -- node[w,above right] {5} (T);
\draw[->] (c) -- node[w,below] {3} (d);
\draw[->] (d) -- node[w,below right] {11} (T);
\draw[->] (a.-50) -- node[w,below left] {1} (d.150);
\draw[->] (d.120) -- node[w,above right] {8} (a.-20);
\end{tikzpicture}
\end{center}
\begin{enumerate}[label=\alph*),nosep]
\item 10
\item 7
\item 16
\item 14
\item None of above.
\end{enumerate}
\end{pastexam}

\begin{pastsoln}
答案：b) 7。

圖上的頂點沒有名字，左邊的是 S、右邊的是 T（題目說從 S 到 T）。 中間四個點叫 $a$（上左）、$b$（上右）、$c$（下左）、$d$（下右）。
邊與容量：$S \to a$ 4、$S \to c$ 6、$a \to b$ 3、$b \to T$ 5、$c \to d$ 3、$d \to T$ 11、$a \to d$ 1、$d \to a$ 8。

\textbf{找流量}（Ford--Fulkerson，資結書第~26.3~節）：
\begin{enumerate}[nosep]
\item $S \to a \to b \to T$，瓶頸 $\min(4, 3, 5) = 3$。
\item $S \to c \to d \to T$，瓶頸 $\min(6, 3, 11) = 3$。
\item $S \to a \to d \to T$，$S \to a$ 剩 1、$a \to d$ 1、$d \to T$ 剩 8，瓶頸 1。
\end{enumerate}
總流量 $3 + 3 + 1 = 7$。 這時 $S \to a$ 已滿，$S \to c$ 還剩 3，但從 $c$ 出去的 $c \to d$ 已滿，剩餘網路裡從 S 只走得到 $c$，
找不到增廣路徑。

\textbf{最小割驗證}（資結書定理~26.6）：取 $\{S, c\}$ 和 $\{a, b, d, T\}$，從前者到後者的邊是 $S \to a$ 4、$c \to d$ 3，容量和 7。
流量 7 等於這個割的容量，所以 7 就是最大流。 （另一個容量也是 7 的割是 $\{S, a, c\}$：$a \to b$ 3、$a \to d$ 1、$c \to d$ 3；
$d \to a$ 是從另一邊指回來的，不算。）

從割 $\{S, c\}$ 也看得出來：不管中間兩條斜邊的容量和方向怎麼改，流量都不會超過 $4 + 3 = 7$。 選項 a) 10 是 S 流出去的
容量總和 $4 + 6$，但 $c$ 收到的 6 只能經過容量 3 的 $c \to d$ 送出去，送不完。
\end{pastsoln}
